% !TEX root = ../main.tex
%% Corrigés — Chapitre 12 — Combinaisons linéaires et indépendance
%% Tous les résultats ont été vérifiés par programme (calcul exact).

\corrigesChapitre{12}{Combinaisons linéaires et indépendance}

\begin{deuxcolonnes}

\begin{corrige}{1}{Rappels : calculs dans le plan}
\begin{questions}
  \item $\vec u+\vec v\,\res{(2\pv6)}$ ; $2\vec u-\vec w\,\res{(1\pv6)}$ ;\\ $\vec u+\vec v+\vec w\,\res{(7\pv4)}$ ;
    $3\vec u-2\vec v+\vec w\,\res{(16\pv-4)}$.
  \item $\vect{AB}\,\res{(5\pv-4)}$.
  \item $2\times7-3\times5=\res{-1}\neq0$ : vecteurs non colinéaires ;\\ $6\times6-4\times9=\res{0}$ : vecteurs colinéaires.
  \item $6m-3\times2=0$, soit $\res{m=1}$.
\end{questions}
\end{corrige}

\begin{corrige}{2}{Colinéarité et alignement}
\begin{questions}
  \item $\vec u=2\vec v$ : $\res{\text{colinéaires}}$. $\det(\vec a,\vec b)=-15+12=-3\neq0$ : $\res{\text{non colinéaires}}$.
  \item $4k-24=0$, soit $\res{k=6}$.
  \item $\vect{AB}\,(3\pv3)$ et $\vect{AC}\,(6\pv6)=2\vect{AB}$ : $\res{\text{alignés}}$.
  \item $\det(\vec u,\vec v)=12-12=0$ : $\res{\text{pas une base}}$ ($\vec v=-2\vec u$). $\det(\vec m,\vec n)=35-6=29\neq0$ :
    $\res{\text{base}}$.
\end{questions}
\end{corrige}

\begin{corrige}{3}{Décomposer un vecteur}
\methode{écrire $\vec w=a\vec u+b\vec v$ et résoudre le système}
\begin{questions}
  \item $\det(\vec u,\vec v)=-1-6=-7\neq0$ : $\vec u$ et $\vec v$ forment une base.
  \item $a+3b=7$ et $2a-b=1$ donnent\\ $\res{\vec w=\tfrac{10}{7}\vec u+\tfrac{13}{7}\vec v}$ ;\\
    $a+3b=5$ et $2a-b=8$ donnent\\ $\res{\vec a=\tfrac{29}{7}\vec u+\tfrac27\vec v}$.
  \item Oui : $\res{\vec b=2\vec u}$.
  \item $a\vec u+(5-a)\vec v$ a pour coordonnées $\res{(15-2a\pv3a-5)}$.
\end{questions}
\end{corrige}

\begin{corrige}{4}{Changer de base dans le plan}
\begin{questions}
  \item $\res{\vec u=2\vec{e_1}+3\vec{e_2}}$.
  \item $\det(\vec u,\vec v)=-2-3=-5\neq0$ : c'est une base.
  \item $2a+b=1$ et $3a-b=0$ : $\res{\vec{e_1}=\tfrac15\vec u+\tfrac35\vec v}$.\\
    $2a+b=0$ et $3a-b=1$ : $\res{\vec{e_2}=\tfrac15\vec u-\tfrac25\vec v}$.
  \item $2a+b=5$ et $3a-b=4$ : $\res{a=\tfrac95}$, $\res{b=\tfrac75}$. Avec la question 3 :\\
    $5\vec{e_1}+4\vec{e_2}=\left(1+\frac45\right)\vec u+\left(3-\frac85\right)\vec v$.
\end{questions}
\end{corrige}

\begin{corrige}{5}{Une famille liée}
\begin{questions}
  \item $\det(\vec{u_1},\vec{u_2})=-1-2=-3\neq0$ : indépendants.
  \item $a+2b=5$ et $a-b=1$ : $\vec{u_3}=\frac73\vec{u_1}+\frac43\vec{u_2}$, d'où $\res{7\vec{u_1}+4\vec{u_2}-3\vec{u_3}=\vec0}$.
  \item Si deux d'entre eux sont colinéaires, la famille est liée. Sinon, ils forment une base du plan, et le troisième s'exprime en
    fonction des deux autres.
  \item $a+2b=0$ et $a-b=5$ : $\res{\vec w=\tfrac{10}{3}\vec{u_1}-\tfrac53\vec{u_2}}$.
\end{questions}
\end{corrige}

\begin{corrige}{6}{Applications géométriques}
\begin{questions}
  \item $\res{M\left(\tfrac72\pv2\right)}$.
  \item $\vect{AB}\,(3\pv2)=\vect{CD}$ : $\res{(AB)\parallel(CD)}$.
  \item $ABCE$ parallélogramme $\iff\vect{AB}=\vect{EC}$, d'où $E=C-\vect{AB}$ : $\res{E(-2\pv2)}$.
  \item $\vect{AC}\,(-1\pv3)$ et $\vect{AD}\,(2\pv5)=\vect{AB}+\vect{AC}$ : $ABDC$ est un $\res{\text{parallélogramme}}$.
  \item $\det(\vect{AB},\vect{AD})=15-4=11\neq0$ : $\res{\text{non alignés}}$.
\end{questions}
\end{corrige}

\begin{corrige}{7}{Équations vectorielles}
\begin{questions}
  \item $2x-y=3$ et $x+3y=7$ : $\res{x=\tfrac{16}{7}}$, $\res{y=\tfrac{11}{7}}$.
  \item $2x-y=1$, $x+3y=5$ : $\res{\vec a=\tfrac87\vec u+\tfrac97\vec v}$ (toujours possible, car $(\vec u,\vec v)$ est une base).
  \item $b=2a$ et $a+3b=10$ : $\res{a=\tfrac{10}{7}}$, $\res{b=\tfrac{20}{7}}$.
  \item De $7\vec w=16\vec u+11\vec v$ :\\ $\res{\vec v=-\tfrac{16}{11}\vec u+\tfrac{7}{11}\vec w}$.
  \item $\vect{MN}\,(-3\pv2)$, $\vect{MP}\,(6\pv-4)=-2\vect{MN}$ : $\res{\text{alignés}}$.
\end{questions}
\end{corrige}

\begin{corrige}{8}{Vecteurs de l'espace}
\begin{questions}
  \item $\vec w=a\vec u+b\vec v$ : $a=2$, $b=1$, et la cote $2a+b=5$ convient : $\res{\vec w=2\vec u+\vec v}$. La famille est
    $\res{\text{liée}}$.
  \item $a\vec u+b\vec v+c\vec t=\vec0$ donne $a=0$, $b=0$, puis $2a+b+c=0$, soit $c=0$ : la famille est libre, c'est une
    $\res{\text{base}}$.
  \item $\vect{OA}=\vec u$, $\vect{OB}=\vec u+\vec v$. Si $\vect{OC}=a\vect{OA}+b\vect{OB}$ : $b=2$, $a=1$, mais la cote vaut
    $2a+3b=8\neq7$. Les points ne sont $\res{\text{pas coplanaires}}$.
\end{questions}
\end{corrige}

\begin{corrige}{9}{Coplanarité et base dans l'espace}
\begin{questions}
  \item $\vec v=k\vec u$ donnerait $k=2$ puis $k=\frac12$ : non colinéaires. Et $\res{\vec w=\vec u+\vec v}$.
  \item Oui, ils sont $\res{\text{coplanaires}}$ : la famille est liée, ce n'est pas une base.
  \item $a\vec u+b\vec v+c\vec t=\vec0$ : $2a+b=0$ et $a+3b=0$ donnent $a=b=0$, puis $c=0$ : $\res{\text{base}}$.
  \item $\vect{OA}=\vec u$, $\vect{OB}=\vec w=\vec u+\vec v$ et $\vect{OC}=\vec u+2\vec v$ : les trois vecteurs sont combinaisons de
    $\vec u$ et $\vec v$, les points sont $\res{\text{coplanaires}}$.
\end{questions}
\end{corrige}

\end{deuxcolonnes}

%% ---------------------------------------------------------------------
\corrigesSouples
\begin{corrige}{10}{Coordonnées dans une base non orthogonale}
\begin{questions}
  \item Si $a\vec u+b\vec v+c\vec w=\vec0$, alors $a+c=0$, $a+b=0$ et $b+c=0$. En additionnant : $2(a+b+c)=0$, donc $a+b+c=0$, puis
    $b=0$, $c=0$ et $a=0$. La famille est libre : c'est une $\res{\text{base}}$.
  \item Pour $\vec s$ : $a+c=4$, $a+b=3$, $b+c=1$. La somme donne $a+b+c=4$, d'où $b=0$, $c=1$, $a=3$ : $\res{\vec s=3\vec u+\vec w}$,
    de coordonnées $(3\pv0\pv1)$. Pour $\vec r$ : $a+b+c=\frac32$ et $a=b=c=\frac12$ : $\res{\vec r=\tfrac12(\vec u+\vec v+\vec w)}$.
  \item $\vect{OM}=2\vec u-\vec v+3\vec w$ a pour coordonnées $(2+3\pv2-1\pv-1+3)$ : $\res{M(5\pv1\pv2)}$ dans
    $(O\pv\vec i,\vec j,\vec k)$.
  \item $\vect{OM}=X\vec u+Y\vec v+Z\vec w$ donne $\res{x=X+Z}$, $\res{y=X+Y}$ et $\res{z=Y+Z}$.
\end{questions}
\end{corrige}

\begin{corrige}{11}{Une famille dépendant d'un paramètre}
\begin{questions}
  \item Pour $m=1$, $\vec u=\vec v=\vec w=(1\pv1\pv1)$ : la famille est $\res{\text{liée}}$.
  \item Pour $m=-2$, $\vec u+\vec v+\vec w=(1-2+1\pv-2+1+1\pv1+1-2)=\vec0$ : relation à coefficients non nuls, la famille est
    $\res{\text{liée}}$.
  \item Le système s'écrit $a+mb+c=0$, $ma+b+c=0$ et $a+b+mc=0$. La différence des deux premières équations donne
    $(1-m)(a-b)=0$, celle des deux dernières $(m-1)(a-c)=0$. Comme $m\neq1$, $a=b=c$, et la première équation devient
    $a(1+m+1)=0$, soit $\res{(m+2)a=0}$.
  \item Si $m\neq1$ et $m\neq-2$, on obtient $a=0$, donc $a=b=c=0$ : la famille est libre. C'est une base
    $\res{\text{si et seulement si } m\notin\{1\pv-2\}}$ (les cas $m=1$ et $m=-2$ ont été traités aux questions 1 et 2).
  \item Pour $m=0$ : $\vec u\,(1\pv0\pv1)$, $\vec v\,(0\pv1\pv1)$, $\vec w\,(1\pv1\pv0)$ et $\vec u+\vec v+\vec w=(2\pv2\pv2)$ :
    $\res{\vec s=\vec u+\vec v+\vec w}$, de coordonnées $(1\pv1\pv1)$ dans cette base.
\end{questions}
\end{corrige}

\begin{corrige}{12}{Centre de gravité d'un tétraèdre}
\begin{questions}
  \item $A(0\pv0\pv0)$, $B(1\pv0\pv0)$, $C(0\pv1\pv0)$, $D(0\pv0\pv1)$. La relation $\vect{GB}+\vect{GC}+\vect{GD}=\vec0$ s'écrit
    $3\vect{AG}=\vect{AB}+\vect{AC}+\vect{AD}$ (relation de Chasles), d'où $\res{G\left(\tfrac13\pv\tfrac13\pv\tfrac13\right)}$.
  \item De même, $4\vect{AO}=\vect{AB}+\vect{AC}+\vect{AD}$ donne $O\left(\frac14\pv\frac14\pv\frac14\right)$. Alors
    $\vect{AO}=\frac34\vect{AG}$ : $A$, $O$, $G$ sont alignés et $O$ est situé $\res{\text{aux trois quarts de }[AG]}$ en partant de $A$.
  \item $I\left(\frac12\pv0\pv0\right)$ et $J\left(0\pv\frac12\pv\frac12\right)$ : le milieu de $[IJ]$ a pour coordonnées
    $\left(\frac14\pv\frac14\pv\frac14\right)$, c'est $\res{O}$.
  \item Le même calcul pour les milieux de $[AC]$ et $[BD]$, puis de $[AD]$ et $[BC]$, donne encore le point $O$ (les quatre sommets
    jouent des rôles symétriques dans la définition de $O$). Les trois segments sont $\res{\text{concourants en }O}$, centre de gravité
    du tétraèdre.
\end{questions}
\end{corrige}
